\section{A Retracted Finding: The Discrete GPU Was Never Broken}
\label{sec:egl}

For three working sessions the recorded conclusion was that the discrete GPU
could not do Vulkan on this system at all---that its driver was, in the words
written down at the time, dead. That conclusion was wrong, it was wrong for a
reason that generalises, and correcting it is what made everything in
Sections~\ref{sec:deadlock} through~\ref{sec:fix} possible.

\subsection{The symptom}

Any attempt to create a Vulkan instance restricted to the NVIDIA driver failed
at the earliest possible point:

\begin{lstlisting}
Could not get 'vkCreateInstance' via
 'vk_icdGetInstanceProcAddr' for ICD libGLX_nvidia.so.0
\end{lstlisting}

The loader could open the driver library but could not obtain a single entry
point from it. Every global function resolved to null. This is not a subtle
failure; it looks exactly like a broken or mismatched driver.

\subsection{Isolating it below the compatibility layer}

The first useful decision was to stop testing through Wine. A twenty-line C
program was written that does nothing but \verb|dlopen| an ICD library, call its
negotiation entry point across every interface version the loader supports, and
attempt one instance creation.

\begin{table}[htbp]
\caption{Direct ICD probe, outside Wine and outside the Vulkan loader}
\label{tab:icd}
\centering
\footnotesize
\begin{tabularx}{\columnwidth}{@{}L{3.4cm}cc@{}}
\toprule
\textbf{Probe} & \textbf{Mesa pin} & \textbf{NVIDIA present} \\
\midrule
\verb|icdNegotiate| v1--7 & $-3$ & $0$ \\
\verb|vkCreateInstance| & $-3$ & $0$ \\
Resulting instance handle & \verb|(nil)| & \verb|0x557e90579430| \\
\bottomrule
\end{tabularx}
\end{table}

The difference between the two columns is a single environment variable. That is
the whole finding: the driver is fine, and something in the environment is
disabling it.

\subsection{Root cause: an environment variable that replaces rather than adds}

The system pins its EGL vendor selection through
\verb|__EGL_VENDOR_LIBRARY_FILENAMES|, set to Mesa's \verb|50_mesa.json| so that
ordinary desktop applications do not incidentally wake the discrete GPU.

\textbf{That variable is a replacement, not an addition.} It does not extend the
vendor search path; it \emph{is} the vendor search path. With only Mesa's entry
visible, libglvnd~\cite{glvnd} has no NVIDIA vendor at all.

The non-obvious part is why this breaks \emph{Vulkan}, which is a separate API
with its own loader and its own ICD list. NVIDIA's Vulkan implementation
initialises through the same vendor-dispatch machinery as its EGL
implementation. With the vendor absent, that initialisation fails silently and
the ICD returns a table of null pointers rather than an error.

The confirming observation is from \verb|strace|: in the failing configuration
the driver library \textbf{never opens \texttt{/dev/nvidiactl}}. It is not being
refused by the kernel; it never asks. This also independently eliminates the
access gate of Section~\ref{sec:system}, which was again the first suspect.

\subsection{Fix}

Whenever the NVIDIA Vulkan driver is required, the EGL vendor list must name
\emph{both} vendors---Mesa's entry and \verb|60_nvidia.json|. The rule to carry
forward:

\begin{quote}
\textbf{A pin that names one vendor is a policy statement, and every consumer of
that policy inherits it, including ones in a different API.} Any launcher that
needs the pinned-out vendor must restore it explicitly rather than assume it can
add to a list.
\end{quote}

After the correction, \verb|vulkaninfo| reports the discrete GPU exactly as it
should: \textbf{NVIDIA GeForce RTX 4050 Laptop GPU}, device type
\verb|DISCRETE_GPU|, Vulkan 1.4.329, driver 595.71.05.

\subsection{A second, independent fault hiding behind the first}

Correcting the vendor list revealed that a second defect had been masked by it.
The project's Wine launcher, when asked for the discrete GPU, was performing a
bare \verb|exec wine| with no ICD selection at all---so it silently used the
integrated Radeon while the desktop notification it emitted said NVIDIA.

Two independent faults, one of which lies to the user about which hardware is in
use, are a good illustration of why the ground-truth rule from
Section~\ref{sec:graphics1} matters. The failing configuration was
indistinguishable from the working one by every self-report available, and
distinguishable instantly by \verb|nvidia-smi --query-compute-apps|.

\subsection{The methodological lesson}

The recorded claim was ``NVIDIA Vulkan is dead on this machine.'' The evidence
for it was a test that had been run under a system-wide policy the test did not
know about.

\begin{quote}
\textbf{When an instrument reports that a major component is entirely
non-functional, suspect the instrument first.} Total failure is rare; misconfigured
measurement is common. The cheapest discriminator is to reproduce the failure at
the lowest possible layer---here, forty lines of C with no loader, no
compatibility layer, and no launcher script between the test and the driver.
\end{quote}

A corollary specific to this class of bug: a bare test of the form ``set the ICD
list and see whether it works'' is \emph{not} a test of the driver, because the
ICD list is not the only list involved. It is a test of one variable in an
environment with at least two.

With the discrete GPU restored to service, the game could finally be run on it.
It then produced the failure that defines the rest of this paper.
